\begin{proposition}[DP36: unequal weights restore the kernel quantization]\label{DP36}
Let $p$ be odd, $h\ge2$, $q=p^h$, $d=h-1$, and $R=W(k)$ with $\F_q\subset k$. Use the actual pair $G_h,H=G_h^D$ of DP32, with logarithms $L_0,\mathcal L$ and inverses $E_0,\mathcal E$. Put
\[
 s=(p-1)^{-1},\quad\rho=(q-1)^{-1},\quad r_0=\rho,\quad r_i=s-p^i\rho\ (1\le i\le d).
\]
Suppose $\beta(x,\mathbf y)\in R[1/p][[x,\mathbf y]]$ vanishes on either axis, has no bidegree-$(1,1)$ monomial, and satisfies
\[
 \beta(\zeta x,\zeta^{-p}y_1,\ldots,\zeta^{-p^d}y_d)=\beta(x,\mathbf y)
\]
for a primitive $(q-1)$-st root $\zeta$. Suppose its two exponential addition coboundaries
\[
 C_X=\exp\bigl(\delta_x\beta(L_0(S),L_0(S'),\mathcal L(\mathbf T))\bigr),\quad
 C_Y=\exp\bigl(\delta_{\mathbf y}\beta(L_0(S),\mathcal L(\mathbf T),\mathcal L(\mathbf T'))\bigr)
\]
are integral units with constant coefficient one. Here $\delta_{\mathbf y}$ uses addition of the whole vector in logarithm coordinates. For any integral series $U$ in the group coordinates and $\sigma\in\{1,-1\}$, set
\[
 \Theta(x,\mathbf y)=U(E_0(x),\mathcal E(\mathbf y))\exp\bigl((\sigma/2)\beta(x,\mathbf y)\bigr).
\]
Then
\begin{equation}\label{eq:true-pair-weight}
 v_p(\mathbf a![x^{a_0}\mathbf y^{\mathbf a'}]\Theta)
 \ge (s-p\rho)a_0+\sum_{i=1}^d\left(\frac{p^{i+1}}{q-1}-1\right)a_i
       -\frac{\sum_{i=0}^d s_p(a_i)}{p-1}.
\end{equation}
On the specific path $x=t$, $y_d=\epsilon t$, $y_i=0$ for $i<d$, with $\epsilon\in R^\times$, this gives
\begin{equation}\label{eq:true-pair-path}
 v_p(n![t^n]\Theta(t,0,\ldots,0,\epsilon t))
 \ge\rho n-\frac{s_p(n)}{p-1}.
\end{equation}
These are sufficient estimates, with no asserted optimality or geometric realization of the hypotheses or path.
\end{proposition}
\begin{proof}
If $x^a\mathbf y^{\mathbf b}$ is invariant, $a-\sum_i p^ib_i$ is divisible by $q-1$. Its weighted degree is
\[
 r_0a+\sum_i r_ib_i
 =\rho\left(a-\sum_i p^ib_i\right)+s|\mathbf b|\in s\Z.
\]
All $r_i$ are positive. Every nonconstant invariant monomial therefore has weighted degree at least $s$. The same calculation applies with duplicated $x$-variables or duplicated $\mathbf y$-vectors. Integral coefficient valuations are nonnegative, so the lower bound remains valid. The actual logarithms, exponentials and group laws are equivariant for these characters by their explicit formulas in DP32. Thus the integral series $C_X-1,C_Y-1$ have weighted valuation at least $s$.

The inverse-logarithm estimates give weight at least $r_0$ to $E_0$ and at least $r_i$ to $\mathcal E_i$; substituting them preserves this bound. Taking the formal logarithm consequently gives
\[
 w_{\mathbf r}(\delta_x\beta),\ w_{\mathbf r}(\delta_{\mathbf y}\beta)\ge s.
\]
Indeed its $m$-th term has weight at least $ms-v_p(m)\ge s$, since $v_p(m)\le(m-1)/(p-1)$.

For a monomial $x^a\mathbf y^{\mathbf b}$ of $\beta$, both side degrees are positive and not both one. If $a\ge2$, choose an interior binomial coefficient of valuation at most one as in DP29 and read the corresponding coefficient of $\delta_x\beta$. If $a=1$ and $|\mathbf b|\ge2$, either two components of $\mathbf b$ are nonzero, in which case a split sending one whole component to the first vector has coefficient one, or a single component has exponent at least two, in which case use the same binomial choice in that component. Recombining the split exponents recovers the original monomial uniquely, so there is no contamination by another coefficient. It follows that $w_{\mathbf r}(\beta)\ge s-1$.

Each monomial of $\beta$ has positive weighted degree at least $s$. Increasing all weights from $r_i$ to $pr_i$ therefore adds at least $(p-1)s=1$ to every coefficient weight. Thus $w_{p\mathbf r}(\beta)\ge s$, and its exponential with multiplier $\sigma/2$ has nonnegative weight by $v_p(m!)\le ms$. The integral $U$ factor has nonnegative weight as well; increasing the positive weights preserves the established inverse bounds. Hence the ordinary coefficient is at least $-p\sum_i r_ia_i$ in valuation. Adding Legendre's exact factorial identity proves \eqref{eq:true-pair-weight}.

For the stated path assign $t$ weight $pr_d$. Since $r_d\ge r_0$, both nonzero substitutions have at least their required weights. Setting the other coordinates to zero preserves the bound. The resulting factorial coefficient is bounded by $(s-pr_d)n-s_p(n)/(p-1)$. The identity $s-pr_d=\rho$ proves \eqref{eq:true-pair-path}.
\end{proof}
